\chapter{Punkty osobliwe i geometria lokalna krzywych algebraicznych}
\label{ch:osobliwosci}

\section{Dlaczego obraz krzywej może być mylący?}

Prosta, parabola i okrąg mają w zwykłych punktach jeden wyraźny
kierunek styczny. Natomiast krzywe $y^2=x^3$ oraz
$y^2=x^2(x+1)$ mają w początku układu współrzędnych
zachowanie zupełnie odmienne od prostej: pierwsza ma
\emph{ostrze}, a druga \emph{węzeł} z dwiema gałęziami.
Sama informacja, że jest to zbiór zerowy jednego wielomianu,
nie mówi jeszcze, jak rozróżnić takie punkty.

Celem tego rozdziału jest zrozumienie trzech poziomów opisu:
\emph{przestrzeni stycznej}, \emph{stożka stycznego} oraz
\emph{lokalnego pierścienia funkcji}. Zobaczymy, że przestrzeń
styczna może mieć zbyt duży wymiar, stożek styczny rozróżnia
kierunki dochodzenia, a pierścień lokalny zachowuje także
ich krotności i inne informacje algebraiczne.

W całym rozdziale $k$ oznacza ciało algebraicznie domknięte;
gdy odwołujemy się do lokalnej intuicji analitycznej,
przyjmujemy $k=\C$. W przykładzie węzła i w obliczeniach
z dwiema odmiennymi stycznymi dodatkowo zakładamy
$\operatorname{char}k\ne2$.

\begin{center}
\begin{tikzpicture}[>=Stealth,scale=.83]
  \begin{scope}[shift={(-4,0)}]
   \draw[->,gray!70] (-.4,0)--(2.7,0) node[right] {$x$};
   \draw[->,gray!70] (0,-1.65)--(0,1.85) node[above] {$y$};
   \draw[blue!75!black,very thick,domain=-1.25:1.25,samples=120]
    plot ({\x*\x},{\x*\x*\x});
   \fill (0,0) circle(2pt); \node at (1.5,-1.55) {ostrze $y^2=x^3$};
  \end{scope}
  \begin{scope}[shift={(2.3,0)}]
   \draw[->,gray!70] (-1.7,0)--(2.8,0) node[right] {$x$};
   \draw[->,gray!70] (0,-1.7)--(0,1.9) node[above] {$y$};
   \draw[red!75!black,very thick,domain=-1.7:1.7,samples=145]
    plot ({\x*\x-1},{\x*(\x*\x-1)});
   \fill (0,0) circle(2pt); \node at (.7,-2.05) {węzeł $y^2=x^2(x+1)$};
  \end{scope}
\end{tikzpicture}

\small Rysunek 3.1. Dwie krzywe o osobliwym początku, lecz o różnej
liczbie lokalnych gałęzi. Na prawym rysunku parametry $t=1$ i $t=-1$
dają ten sam punkt $(0,0)$.
\end{center}

\section{Pierścień lokalny: funkcje w pobliżu punktu}

\begin{definicja}[Pierścień lokalny w punkcie]
\label{def:local-ring-point}
Niech $X\subseteq k^n$ będzie niepustym zbiorem algebraicznym,
$A=k[X]$ jego pierścieniem współrzędnych, a $p\in X$.
Przez $\mathfrak m_p\subseteq A$ oznaczamy ideał klas
funkcji wielomianowych znikających w $p$.
\emph{Pierścieniem lokalnym} $X$ w punkcie $p$ nazywamy
\[
 \mathcal O_{X,p}=A_{\mathfrak m_p}
 =\left\{\frac fg:\ f,g\in A,\ g(p)\ne0\right\}/\!\sim,
\]
gdzie $f/g=f'/g'$ wtedy i tylko wtedy, gdy dla pewnego
$h\in A$ z $h(p)\ne0$ mamy $h(fg'-f'g)=0$ w $A$.
\end{definicja}

Warunek dodatkowego czynnika $h$ w relacji równości jest
istotny w pierścieniach z dzielnikami zera. Jeśli $A$
jest dziedziną, wystarczy zwykła równość $fg'=f'g$.
Funkcja $f/g$ ma sens na pewnym otoczeniu punktu $p$,
ponieważ $g(p)\ne0$, a zbiór $D(g)$ jest otwarty
w topologii Zariskiego.

\begin{twierdzenie}[Jedyny ideał maksymalny]
\label{tw:local-unique-max}
Dla każdego algebraicznego $X\subseteq k^n$ i $p\in X$
pierścień $\mathcal O_{X,p}$ ma dokładnie jeden ideał
maksymalny, równy
$\mathfrak m_p\mathcal O_{X,p}=\{f/g:f(p)=0,\ g(p)\ne0\}$.
Element $f/g$ jest odwracalny dokładnie wtedy, gdy $f(p)\ne0$.
\end{twierdzenie}
\begin{proof}
Podstawienie w $p$ jest dobrze określonym homomorfizmem
$\operatorname{ev}_p:\mathcal O_{X,p}\to k$,
$f/g\mapsto f(p)/g(p)$: jeśli ułamki są równe,
po podstawieniu w $p$ ich iloczyny na krzyż są równe.
Homomorfizm jest surjektywny, bo przyjmuje każdą stałą.
Jego jądro to dokładnie ułamki o liczniku znikającym w $p$,
a więc jest to ideał maksymalny.

Jeżeli $f(p)\ne0$, to $g/f$ jest poprawnym ułamkiem
lokalnym i jest odwrotnością $f/g$. Jeśli $f(p)=0$,
to $\operatorname{ev}_p(f/g)=0$, a więc element ten nie
może być odwracalny: homomorfizm zachowujący jedynkę
przenosi jednostki na jednostki.
W każdym pierścieniu ideał maksymalny nie zawiera jednostek.
Dowolny ideał właściwy $J$ może więc składać się tylko
z niejednostek i musi zawierać się w jądrze
$\operatorname{ev}_p$. Jedyny ideał maksymalny jest zatem
właśnie tym jądrem.
\end{proof}

\begin{przyklad}[Regularna funkcja na otoczeniu, ale nie na całej krzywej]
Na hiperboli $H=\V_k(xy-1)$ mamy
$k[H]\cong k[x,x^{-1}]$; element $1/x$ jest nawet
globalną funkcją regularną, bo w pierścieniu współrzędnych
$1/x=y$. Natomiast na prostej $\A^1(k)$ funkcja $1/x$
nie należy do $k[x]$, lecz należy do lokalnego pierścienia
w każdym punkcie $p\ne0$. Nie należy do
$\mathcal O_{\A^1,0}$, ponieważ mianownik $x$ znika w $0$.
\end{przyklad}

\begin{przyklad}[Lokalne równania węzła]
Dla $X=\V_k(y^2-x^2(x+1))$ i $p=(0,0)$
mamy $\mathcal O_{X,p}=(k[x,y]/(y^2-x^2(x+1)))_{(x,y)}$.
Element $1+x$ jest jednostką tego pierścienia, ponieważ
$(1+x)(p)=1$. Nie znaczy to, że $1+x$ jest jednostką
całego pierścienia $k[X]$.
\end{przyklad}

\section{Przestrzeń kostyczna i przestrzeń styczna Zariskiego}

\begin{definicja}[Przestrzeń kostyczna i styczna]
\label{def:tangent-zariski}
Niech $X\subseteq k^n$ będzie zbiorem algebraicznym,
$p\in X$, a $\mathfrak m$ jedynym ideałem maksymalnym
$\mathcal O_{X,p}$. Iloraz $\mathfrak m/\mathfrak m^2$
traktujemy jako przestrzeń wektorową nad
$\mathcal O_{X,p}/\mathfrak m\cong k$; nazywa się ją
\emph{przestrzenią kostyczną} w $p$.
\emph{Przestrzeń styczna Zariskiego} to
\[
 T_pX:=\operatorname{Hom}_k(\mathfrak m/\mathfrak m^2,k).
\]
\end{definicja}

Dlaczego modulo $\mathfrak m^2$? Iloczyny dwóch funkcji
znikających w $p$ są \emph{drugiego rzędu małości}.
W zwykłym rachunku różniczkowym pochodna iloczynu
takich funkcji w $p$ jest zerowa. Gdy więc zostawimy
jedynie klasy modulo $\mathfrak m^2$, zachowamy
informację o wyrazach pierwszego rzędu.

\begin{twierdzenie}[Jawny wzór na przestrzeń styczną]
\label{tw:zariski-tangent-formula}
Niech $k$ będzie ciałem, $X=\V_k(I)\subseteq k^n$
zbiorem algebraicznym, $p=(a_1,\dots,a_n)\in X$,
a $I(X)$ jego ideałem znikania. Wtedy istnieje
naturalna identyfikacja
\[
 T_pX\cong\left\{v\in k^n:
 \sum_{j=1}^n\frac{\partial f}{\partial x_j}(p)v_j=0
 \quad\text{dla każdego }f\in I(X)\right\}.
\]
Jeżeli $I(X)=(f_1,\dots,f_r)$, to wystarczy sprawdzać
te równania dla generatorów $f_i$.
\end{twierdzenie}
\begin{proof}
Przesuńmy współrzędne tak, aby $p$ stał się początkiem.
Niech $R=k[x_1,\dots,x_n]$ i
$\mathfrak n=(x_1-a_1,\dots,x_n-a_n)$.
Każdy wielomian z $\mathfrak n$ można zapisać jako
sumę wyrazów liniowych w $x_j-a_j$ i wyrazów należących
do $\mathfrak n^2$. W rezultacie
$\mathfrak n/\mathfrak n^2$ ma bazę klas
$x_1-a_1,\dots,x_n-a_n$ i jest izomorficzny z $k^n$.

W pierścieniu $A=k[X]=R/I(X)$ ideał punktu jest obrazem
$\mathfrak n$. Lokalizacja w $p$ nie zmienia ilorazu
$\mathfrak m/\mathfrak m^2$: każdy mianownik nieznikający
w $p$ ma w tym ilorazie tę samą wartość co jego stała
wartość w $p$, ponieważ różnica należy do $\mathfrak m$,
a jej iloczyn przez element $\mathfrak m$ należy do
$\mathfrak m^2$. Zatem przestrzeń kostyczna na $X$
jest ilorazem $\mathfrak n/\mathfrak n^2$ przez obrazy
wielomianów $f\in I(X)$, których wyrazy stałe są równe zero.

Wyraz liniowy rozwinięcia $f$ w $p$ ma postać
$\sum_j(\partial f/\partial x_j)(p)(x_j-a_j)$.
Przestrzeń dualna ilorazu jest więc anihilatorem
wszystkich tych form liniowych w $(k^n)^*$.
To dokładnie podany zbiór wektorów $v$.
Jeśli $f=\sum_i h_if_i$ oraz $f_i(p)=0$, to
reguła Leibniza daje
$df(p)=\sum_i h_i(p)\,df_i(p)$, więc wystarcza
sprawdzenie generatorów ideału.
\end{proof}

\begin{uwaga}[Dlaczego nie wystarczą generatory dowolnego ideału równań?]
Jeżeli $I$ nie jest radykalny, zbiór punktów $\V_k(I)$
może mieć większy ideał znikania $I(X)=\sqrt I$.
Dla zbioru $X=\{0\}\subseteq\A^1$ wyznaczonego
równaniem $x^2=0$ pochodna podanego równania w $0$
wynosi zero i błędnie sugeruje jednowymiarową przestrzeń
styczną. Jednak jako \emph{zredukowany zbiór punktów}
$X$ ma ideał znikania $(x)$ i $T_0X=0$.
Jeśli świadomie badamy niezredukowany obiekt
$\operatorname{Spec}k[x]/(x^2)$, jego przestrzeń
styczna ma wymiar $1$. Rozróżnienie to będzie kluczowe
w geometrii schematów, choć nie rozwijamy tu ich teorii.
\end{uwaga}

\begin{przyklad}[Przestrzeń styczna w punkcie gładkiej paraboli]
Niech $X=\V_k(y-x^2)$ i $p=(a,a^2)$.
Ponieważ $f_x=-2x$ oraz $f_y=1$, mamy
$T_pX=\{(u,v)\in k^2:-2au+v=0\}$,
czyli jest to prosta o kierunku $(1,2a)$.
Zgadza się to z pochodną parametryzacji $t\mapsto(t,t^2)$.
\end{przyklad}

\begin{przyklad}[Przestrzeń styczna w ostrzu i węźle]
W początku krzywej $y^2=x^3$ pochodne $-3x^2$ i $2y$
są zerowe. Podobnie dla $y^2=x^2(x+1)$
obie pochodne znikają w $(0,0)$. W obu przypadkach
wzór na przestrzeń styczną nie nakłada żadnego warunku
na $(u,v)$, a więc $T_{(0,0)}X=k^2$.
Krzywa ma jednak lokalnie wymiar $1$. Otrzymaliśmy
zatem \emph{za dużą} przestrzeń styczną, co sygnalizuje
osobliwość, ale nie rozróżnia jeszcze dwóch jej typów.
\end{przyklad}

\section{Kryterium Jacobiego dla krzywych płaskich}

\begin{definicja}[Gładki punkt krzywej afinicznej]
\label{def:smooth-affine-curve}
Niech $k$ będzie algebraicznie domkniętym ciałem,
a $X=\V_k(f)\subseteq k^2$ zredukowaną krzywą płaską,
czyli $f\in k[x,y]$ jest niezerowy, niestały i bezkwadratowy.
Punkt $p\in X$ nazywamy \emph{gładkim}, jeżeli
$(f_x(p),f_y(p))\ne(0,0)$; w przeciwnym razie
jest \emph{osobliwy}. Gładkość może być też opisana
warunkiem $\dim_kT_pX=1$.
\end{definicja}

\begin{stwierdzenie}[Równoważność kryteriów dla jednej krzywej]
Dla krzywej z definicji~\ref{def:smooth-affine-curve}
warunki $\dim_k T_pX=1$ i
$(f_x(p),f_y(p))\ne(0,0)$ są równoważne.
W punkcie osobliwym mamy $\dim_kT_pX=2$.
\end{stwierdzenie}
\begin{proof}
Ponieważ $f$ jest bezkwadratowy i $k$ jest algebraicznie
domknięte, silne Nullstellensatz daje $I(X)=(f)$:
ideał główny $(f)$ jest radykalny w UFD, gdy $f$ nie ma
powtórzonych czynników nierozkładalnych.
Z twierdzenia~\ref{tw:zariski-tangent-formula} dostajemy
$T_pX=\{(u,v):f_x(p)u+f_y(p)v=0\}$.
Jest to jądro jednego funkcjonału liniowego $k^2\to k$.
Jeżeli funkcjonał jest niezerowy, jego jądro ma wymiar $1$;
jeżeli jest zerowy, jądro jest całym $k^2$ i ma wymiar $2$.
\end{proof}

\begin{uwaga}[Jak działa twierdzenie o funkcji uwikłanej nad $\C$?]
Jeżeli $k=\C$ i $f_y(p)\ne0$, holomorficzne twierdzenie
o funkcji uwikłanej pokazuje, że w dostatecznie małym
otoczeniu $p=(a,b)$ zbiór $f(x,y)=0$ jest wykresem
funkcji holomorficznej $y=\psi(x)$ z $\psi(a)=b$.
Wówczas $f(x,\psi(x))=0$ i po zróżniczkowaniu
$\psi'(a)=-f_x(p)/f_y(p)$. Nie jest to tylko formalna
analogia z analizą: w gładkim punkcie zespolona krzywa
algebraiczna jest lokalnie jednowymiarową rozmaitością
zespoloną. W punkcie osobliwym twierdzenie o funkcji
uwikłanej nie ma zastosowania w takiej postaci.
\end{uwaga}

\section{Stożek styczny: rozróżniamy kierunki i krotności}

\begin{definicja}[Najniższa część jednorodna]
Niech $f\in k[x,y]$ będzie niezerowym wielomianem
spełniającym $f(0,0)=0$. Zapiszmy go jednoznacznie jako
$f=f_m+f_{m+1}+\cdots+f_d$, gdzie każde $f_j$
jest jednorodne stopnia $j$, $f_m\ne0$ i $m\ge1$.
Wielomian $f_m$ nazywamy \emph{formą początkową} $f$
w początku, a liczbę $m$ \emph{rzędem znikania} $f$
w tym punkcie i oznaczamy $\operatorname{ord}_0f=m$.
Dla punktu $p=(a,b)$ przesuwamy współrzędne o $p$.
\end{definicja}

\begin{definicja}[Stożek styczny krzywej]
Niech $X=\V_k(f)\subseteq k^2$ będzie krzywą zredukowaną,
a $p\in X$. Po przesunięciu $p$ do początku wybierzmy
formę początkową $f_m$ jej zredukowanego równania.
\emph{Stożkiem stycznym} krzywej w $p$ na poziomie zbioru
nazywamy $C_pX=\V_k(f_m)\subseteq k^2$.
Wielomian $f_m$, wraz z krotnościami czynników,
nazywamy \emph{równaniem stożka stycznego z krotnościami}.
\end{definicja}

\begin{lemat}[Forma początkowa iloczynu]
Niech $f,g\in k[x,y]$ będą niezerowe i niech
$\operatorname{ord}_0f=m$, $\operatorname{ord}_0g=r$.
Wtedy $\operatorname{ord}_0(fg)=m+r$ i
$(fg)_{m+r}=f_mg_r$.
\end{lemat}
\begin{proof}
W iloczynie $(f_m+f_{m+1}+\cdots)(g_r+g_{r+1}+\cdots)$
nie ma wyrazów stopnia mniejszego niż $m+r$,
a wyraz stopnia $m+r$ jest dokładnie $f_mg_r$.
Ponieważ pierścień wielomianów nad ciałem jest dziedziną,
$f_mg_r\ne0$. Nie następuje więc skasowanie całej
najniższej części.
\end{proof}

Z lematu wynika, że jeśli ideał krzywej jest główny
$I(X)=(f)$, formy początkowe jego elementów generują
ideał $(f_m)$. Dla ideału o wielu generatorach nie wystarcza
jednak branie najniższej części każdego \emph{podanego}
generatora: w kombinacjach mogą znikać wyrazy najniższych
stopni. W ogólnym przypadku definiuje się ideał początkowy
z \emph{wszystkich} elementów ideału.

\begin{przyklad}[Trzy rodzaje lokalnego zachowania]
Dla $f=y-x^2$ najniższa część ma stopień $1$ i wynosi $y$.
Stożek styczny jest prostą $y=0$, zgodną ze zwykłą styczną.
Dla $g=y^2-x^3$ najniższa część to $y^2$:
stożek styczny jako \emph{zbiór} jest prostą $y=0$,
ale z krotnością $2$. Dla
$h=y^2-x^2(x+1)=y^2-x^2-x^3$ najniższa część
wynosi $y^2-x^2=(y-x)(y+x)$, co daje
\emph{dwie różne} styczne $y=x$ i $y=-x$,
gdy $\operatorname{char}k\ne2$.
\end{przyklad}

\begin{center}
\begin{tikzpicture}[scale=.75,>=Stealth]
 \begin{scope}[shift={(-4.15,0)}]
 \draw[->,gray!65] (-1.7,0)--(1.7,0) node[right] {$x$};
 \draw[->,gray!65] (0,-1.45)--(0,1.6) node[above] {$y$};
 \draw[blue!75!black,very thick] (-1.4,0)--(1.4,0);
 \node[below,align=center,font=\small] at (0,-1.65) {ostrze\\podwójna styczna};
 \end{scope}
 \begin{scope}[shift={(2.3,0)}]
 \draw[->,gray!65] (-1.7,0)--(1.7,0) node[right] {$x$};
 \draw[->,gray!65] (0,-1.45)--(0,1.6) node[above] {$y$};
 \draw[red!75!black,very thick] (-1.4,-1.4)--(1.4,1.4);
 \draw[red!75!black,very thick] (-1.4,1.4)--(1.4,-1.4);
 \node[below,align=center,font=\small] at (0,-1.65) {węzeł\\dwie styczne};
 \end{scope}
\end{tikzpicture}

\small Rysunek 3.2. Stożek styczny ostrza i węzła. W obu przypadkach
przestrzeń styczna Zariskiego ma wymiar $2$, ale stożki styczne
rozróżniają lokalne kierunki gałęzi.
\end{center}

\section{Krotność punktu krzywej i styczność przecięć}

\begin{definicja}[Krotność punktu krzywej płaskiej]
Niech $X=\V_k(f)\subseteq k^2$ będzie zredukowaną krzywą,
a $p\in X$. \emph{Krotnością krzywej w punkcie $p$}
nazywamy $\operatorname{mult}_pX=\operatorname{ord}_pf$,
czyli stopień najniższej niezerowej jednorodnej części
przesuniętego wielomianu $f$.
\end{definicja}

\begin{stwierdzenie}[Krotność a gładkość]
Dla zredukowanej płaskiej krzywej $X=\V_k(f)$
i punktu $p\in X$ następujące warunki są równoważne:
$p$ jest gładki, $\operatorname{mult}_pX=1$ oraz
najniższa część jednorodna równania $f$ jest niezerową
formą liniową.
\end{stwierdzenie}
\begin{proof}
Po przesunięciu $p$ do początku część liniowa $f$
ma postać $f_x(p)x+f_y(p)y$. Jest niezerowa dokładnie wtedy,
gdy co najmniej jedna z pochodnych jest niezerowa.
Z definicji to właśnie warunek $\operatorname{ord}_pf=1$.
\end{proof}

\begin{przyklad}[Ostrze i węzeł mają tę samą krotność]
Dla ostrza $y^2=x^3$ i węzła $y^2=x^2(x+1)$
mamy $\operatorname{mult}_{(0,0)}=2$.
Krotność nie rozróżnia tych osobliwości; potrzebny jest
stożek styczny lub normalizacja.
\end{przyklad}

\begin{przyklad}[Styczna do ostrza ma krotność przecięcia trzy]
Niech $C=\V_k(y^2-x^3)$ oraz $L=\V_k(y)$.
W punkcie $(0,0)$ lokalny pierścień przecięcia ma postać
\[
 \bigl(k[x,y]/(y^2-x^3,y)\bigr)_{(x,y)}
 \cong k[x]_{(x)}/(x^3).
\]
Ma on bazę $1,x,x^2$ nad $k$, więc
$I_{(0,0)}(C,L)=3$. Dla prostej $x=0$ analogicznie
otrzymujemy $k[y]_{(y)}/(y^2)$, a więc krotność $2$.
Prosta styczna $y=0$ spotyka ostrze z wyższą krotnością
niż prosta poprzeczna $x=0$.
\end{przyklad}

\section{Normalizacja: rozdzielanie gałęzi}

\begin{definicja}[Element całkowity i domknięcie całkowite]
Niech $A$ będzie dziedziną całkowitości, a
$K=\Frac(A)$ jej ciałem ułamków. Element $u\in K$
jest \emph{całkowity nad $A$}, gdy spełnia moniczne
równanie wielomianowe o współczynnikach z $A$.
Zbiór wszystkich elementów $K$ całkowitych nad $A$
nazywamy \emph{domknięciem całkowitym} $A$ w $K$
i oznaczamy $\overline A$. Dziedzinę $A$ nazywamy
\emph{normalną}, gdy $\overline A=A$.
\end{definicja}

\begin{uwaga}[Dlaczego jest to nowa informacja?]
Pierścień współrzędnych krzywej z ostrzem jest dziedziną,
a krzywa jest nierozkładalna. Mimo to pierścień ten
nie jest normalny. Brak normalności wykrywa więc
zjawisko, którego nie wykrywa sam brak dzielników zera.
\end{uwaga}

\begin{twierdzenie}[Normalizacja krzywej z ostrzem]
\label{tw:normalize-cusp}
Niech $k$ będzie algebraicznie domknięte i
$A=k[x,y]/(y^2-x^3)$. Wtedy $A$ jest dziedziną,
$\Frac(A)\cong k(t)$ przez $x=t^2$, $y=t^3$,
a domknięcie całkowite $A$ w jego ciele ułamków jest
izomorficzne z $k[t]$. Odpowiadające odwzorowanie
krzywych to $\nu:\A^1(k)\to\V_k(y^2-x^3)$,
$t\mapsto(t^2,t^3)$, które jest bijekcją, lecz nie jest
izomorfizmem algebraicznym.
\end{twierdzenie}
\begin{proof}
W poprzednim rozdziale afinicznym pokazaliśmy szczegółowo,
że homomorfizm $k[x,y]\to k[t]$,
$x\mapsto t^2$, $y\mapsto t^3$, ma jądro $(y^2-x^3)$.
Dlatego $A\cong k[t^2,t^3]\subseteq k[t]$ i jest dziedziną.
W ciele ułamków $A$ mamy $t=y/x$, ponieważ
$t^3/t^2=t$ i $x\ne0$ jako element dziedziny.
Zatem $\Frac(A)=k(t)$.

Element $t$ spełnia moniczne równanie $T^2-x=0$
o współczynnikach z $A$ i jest całkowity nad $A$.
Zatem cały $k[t]=A[t]$ jest skończonym modułem nad $A$:
każda jego potęga redukuje się do kombinacji
$1,t$ ze współczynnikami w $A$.
W szczególności każdy element $k[t]$ jest całkowity nad $A$
na mocy lematu~\ref{lem:determinant-integral}.

Odwrotnie, niech $u\in k(t)$ będzie całkowity nad $A$.
Ponieważ $A\subseteq k[t]$, to samo moniczne równanie
pokazuje, że $u$ jest całkowity nad $k[t]$.
Pierścień $k[t]$ jest pierścieniem euklidesowym, a więc
UFD; pierścień UFD jest całkowicie domknięty w swoim
ciele ułamków. Krótko uzasadnijmy ostatni fakt:
jeśli $u=a/b$ jest ułamkiem nieskracalnym i spełnia
$u^m+c_{m-1}u^{m-1}+\cdots+c_0=0$, to po pomnożeniu
przez $b^m$ otrzymujemy
$a^m=-b(c_{m-1}a^{m-1}+\cdots+c_0b^{m-1})$.
Jeżeli $b$ byłby nieodwracalny, pewien jego czynnik
nierozkładalny dzieliłby $a^m$, a więc $a$, wbrew
nieskracalności. Zatem $b$ jest jednostką i $u\in k[t]$.
Otrzymaliśmy $\overline A=k[t]$.

Ostatnie zdanie o bijektywności i braku izomorfizmu
udowodniliśmy w przykładzie~\ref{pr:cusp}.
\end{proof}

\begin{przyklad}[Normalizacja węzła rozdziela dwa punkty]
\label{pr:node-normalization}
Przyjmijmy $\operatorname{char}k\ne2$ i rozważmy
$C=\V_k(y^2-x^2(x+1))$. Odwzorowanie
\[
 \nu:\A^1(k)\to C,\qquad
 t\longmapsto(t^2-1,\ t(t^2-1))
\]
ma wartości w $C$, ponieważ dla $x=t^2-1$ zachodzi
$y^2=t^2x^2=(x+1)x^2$.
Dla $x\ne0$ parametr jest określony jednoznacznie
przez $t=y/x$. W punkcie $x=0$, $y=0$
warunek $t^2-1=0$ daje \emph{dwa} parametry $t=1$ i $t=-1$.
Dla porównania ostrze ma nad osobliwym punktem tylko
jeden parametr $t=0$.

Homomorfizm $k[x,y]/(y^2-x^2(x+1))\to k[t]$
jest różnowartościowy: każdy wielomian po redukcji modulo
równanie krzywej ma postać $A(x)+yB(x)$, a jego obraz
ma postać $A(t^2-1)+t(t^2-1)B(t^2-1)$.
Pierwszy składnik jest parzysty w $t$, drugi nieparzysty;
przy $\operatorname{char}k\ne2$ ich suma może znikać
jako wielomian tylko wtedy, gdy oba są zerowe.
Podstawienie $x=t^2-1$ jest różnowartościowe, więc
$A=B=0$. Otrzymujemy w ten sposób pierścień współrzędnych
jako podpierścień $k[t]$.
Ponieważ $t^2=x+1$ i $t=y/x$ w ciele ułamków,
rozszerzenie do $k[t]$ jest skończone i całkowite,
a ten ostatni pierścień jest normalny. Tym samym
$\nu$ realizuje normalizację krzywej z węzłem.
\end{przyklad}

\section{Rozdmuchanie punktu: zastąpienie punktu kierunkami}

\begin{definicja}[Rozdmuchanie płaszczyzny w początku]
\label{def:blowup-plane}
W przestrzeni $\A^2(k)\times\PP^1(k)$
o współrzędnych $((x,y),[s:t])$ rozważmy zbiór
\[
 B=\{((x,y),[s:t]):xt=ys\}.
\]
Odwzorowanie $\pi:B\to\A^2(k)$,
$\pi((x,y),[s:t])=(x,y)$, nazywamy
\emph{rozdmuchaniem płaszczyzny w początku}.
Włókno $E=\pi^{-1}(0,0)\cong\PP^1(k)$ to
\emph{dzielnik wyjątkowy}.
\end{definicja}

\begin{stwierdzenie}[Co dzieje się z punktami?]
Dla każdego $(x,y)\ne(0,0)$ włókno
$\pi^{-1}(x,y)$ składa się z jednego punktu,
o współrzędnej rzutowej $[s:t]=[x:y]$.
Nad $(0,0)$ całe $\PP^1(k)$ jest włóknem.
Zbiór $B$ pokrywają dwie mapy afiniczne:
przy $s\ne0$ mamy współrzędne $(u,v)$ określone przez
$x=u$, $y=uv$; przy $t\ne0$ mamy
$x=uv$, $y=u$.
\end{stwierdzenie}
\begin{proof}
Jeśli $(x,y)\ne0$, równanie $xt=ys$ oznacza,
że wektory $(x,y)$ i $(s,t)$ są proporcjonalne,
a więc wyznaczają tę samą klasę $[x:y]$.
Dla $(x,y)=0$ równanie jest tożsamością $0=0$,
więc nie ogranicza $[s:t]$.
Gdy $s\ne0$, piszemy $[s:t]=[1:v]$;
z $xt=y$ otrzymujemy $y=xv$. Ustawienie $u=x$
daje zapowiedziane współrzędne. Gdy $t\ne0$, piszemy
$[s:t]=[v:1]$, skąd $x=yv$ i można przyjąć $u=y$.
W obu mapach dzielnik wyjątkowy ma równanie $u=0$.
\end{proof}

\noindent\begin{minipage}{\linewidth}\centering
\begin{tikzpicture}[>=Stealth,scale=.87]
  % Left: image in A^2; right: symbolic local chart of blow-up.
  \draw[->,gray!70] (-2.45,0)--(2.35,0) node[right] {$x$};
  \draw[->,gray!70] (0,-1.85)--(0,1.95) node[above] {$y$};
  \fill (0,0) circle (2.3pt);
  \draw[blue!70!black,thick] (-1.65,-1.2)--(1.65,1.2);
  \draw[red!70!black,thick] (-1.6,1.3)--(1.6,-1.3);
  \node[anchor=north,font=\small] at (0,-2.02) {płaszczyzna $\A^2$};
  % A genuine blow-up map points from the new space to the original plane.
  \draw[<-,thick] (2.65,0)--(4.35,0) node[midway,above] {$\pi$};
  \draw[rounded corners=5pt,thick,gray!70,fill=gray!6]
      (4.55,-1.7) rectangle (8.35,1.65);
  \draw[thick,densely dashed] (6.15,-1.25)--(6.15,1.25);
  \draw[blue!70!black,thick] (4.77,-1.12)--(6.15,-.23);
  \draw[red!70!black,thick] (4.77,.95)--(6.15,.23);
  \fill[blue!70!black] (6.15,-.23) circle (2pt);
  \fill[red!70!black] (6.15,.23) circle (2pt);
  \node[anchor=west,font=\small,fill=white,inner sep=1pt]
      at (6.35,.95) {$E\cong\PP^1$};
  \node[anchor=north,font=\small] at (6.45,-1.83) {lokalny szkic $\widetilde{\A^2}$};
\end{tikzpicture}

\small Rysunek 3.3. Szkic lokalny (nie obraz metryczny): morfizm $\pi$ wysyła cały dzielnik wyjątkowy na początek; rozdmuchanie zastępuje początek
całą prostą rzutową kierunków. Rzeczywiste rozdmuchanie nie jest
po prostu płaską figurą naszkicowaną po prawej.
\end{minipage}

\begin{definicja}[Transformata właściwa krzywej]
Niech $C\subseteq\A^2(k)$ będzie krzywą zawierającą
$(0,0)$. \emph{Transformatą właściwą} $\widetilde C\subseteq B$
nazywamy domknięcie zbioru $\pi^{-1}(C\setminus\{(0,0)\})$
w topologii Zariskiego na $B$.
\end{definicja}

\begin{przyklad}[Rozdmuchanie ostrza]
Dla $C=\V_k(y^2-x^3)$ w mapie $x=u$, $y=uv$
podstawiamy do równania krzywej:
$u^2v^2-u^3=u^2(v^2-u)=0$.
Czynnik $u^2$ opisuje składową wspartą na dzielniku
wyjątkowym i nie należy do równania \emph{transformaty
właściwej}. Aby otrzymać jej równanie, na zbiorze $u\ne0$
dzielimy przez $u^2$, a następnie bierzemy domknięcie.
Otrzymujemy $\widetilde C:\ v^2=u$.

Krzywa $v^2=u$ jest gładka: pochodna równania
$v^2-u$ względem $u$ jest stałą $-1$.
W szczególności osobliwość samej krzywej zniknęła po
jednym rozdmuchaniu. Jednak transformata właściwa
jest styczna do dzielnika wyjątkowego $E:\ u=0$
w ich wspólnym punkcie $(u,v)=(0,0)$, ponieważ
wzór $u=v^2$ ma w początku kierunek $u=0$.
Gdy celem jest uzyskanie przecięć normalnych, potrzebne
są dalsze operacje.
\end{przyklad}

\begin{przyklad}[Rozdmuchanie węzła]
Dla $C=\V_k(y^2-x^2(x+1))$ i $\operatorname{char}k\ne2$
w tej samej mapie podstawienie $x=u$, $y=uv$ daje
$u^2(v^2-u-1)=0$; transformata właściwa spełnia
$v^2=u+1$. Na dzielniku wyjątkowym $u=0$ pojawiają się
dwa punkty $(0,1)$ i $(0,-1)$. Są różne i gładkie,
a styczne transformaty nie pokrywają się tam ze styczną
do dzielnika $u=0$: różniczka
$2v\,dv-du$ ma w każdym z tych punktów niezerowy
współczynnik przy $dv$. Rozdmuchanie rozdzieliło
obie gałęzie węzła.
\end{przyklad}

\section{Wstęp do rodzaju krzywej: algebra spotyka topologię}

Dla gładkich zespolonych krzywych rzutowych geometria
algebraiczna łączy się z topologią powierzchni zwartych.
Krzywa zredukowana, nierozkładalna i gładka nad $\C$
jest zwartą powierzchnią Riemanna, a więc ma rodzaj
$g\in\Z_{\ge0}$, oznaczający liczbę jej uchwytów.
Dla krzywych osobliwych należy rozróżnić rodzaj
arytmetyczny i rodzaj gładkiej normalizacji.

\begin{definicja}[Rodzaj arytmetyczny płaskiej krzywej z ustaloną strukturą]
Niech $k$ będzie ciałem, $S=k[X,Y,Z]$, a $F\in S$
niezerową formą jednorodną stopnia $d\ge1$.
Rozpatrujemy \emph{krzywą określoną przez jednorodny
ideał $(F)$}, czyli nie tylko zbiór punktów
$\V_+(F)$, lecz także stopniowany pierścień $S/(F)$.
Dla $m$ dostatecznie dużego jego funkcja Hilberta ma
postać $\dim_k(S/(F))_m=dm+1-p_a(F)$.
Liczbę $p_a(F)$ nazywamy \emph{rodzajem arytmetycznym}
tej krzywej ze strukturą określoną przez $(F)$.
Dla zredukowanej krzywej płaskiej jej równanie jest
wyznaczone z dokładnością do niezerowego skalara,
więc zapis $p_a(C)$ nie powoduje niejednoznaczności.
\end{definicja}

\begin{przyklad}[Dlaczego sam zbiór punktów nie wystarcza?]
Równania $F=X$ oraz $G=X^3$ wyznaczają ten sam
zbiór punktów $\V_+(F)=\V_+(G)$ w $\PP^2(k)$.
Jednak $S/(X)$ i $S/(X^3)$ nie są izomorficznymi
pierścieniami stopniowanymi: drugi zawiera niezerowe
nilpotenty. Z poniższego twierdzenia wynika
$p_a(F)=0$, natomiast $p_a(G)=1$.
Dlatego rodzaju arytmetycznego struktury \emph{nie wolno}
przypisywać dowolnemu równaniu tylko na podstawie
jego zbioru zer, bez określenia krotności.
\end{przyklad}

\begin{twierdzenie}[Rodzaj arytmetyczny krzywej płaskiej]
\label{tw:arithmetic-genus-plane}
Niech $k$ będzie ciałem, $S=k[X,Y,Z]$, a $F\in S$
niezerową formą stopnia $d\ge1$. Dla każdego $m\ge d$
funkcja Hilberta ilorazu $S/(F)$ jest równa $dm+1-p_a$,
gdzie $p_a=(d-1)(d-2)/2$.
\end{twierdzenie}
\begin{proof}
Mnożenie przez $F$ jest różnowartościowe,
ponieważ pierścień $k[X,Y,Z]$ jest dziedziną.
W dokładnym ciągu stopniowanym
$0\to S(-d)\xrightarrow{\cdot F}S\to S/(F)\to0$
w składowej stopnia $m\ge d$ dostajemy
\[
 \dim_k(S/(F))_m=\binom{m+2}{2}-\binom{m-d+2}{2}
 =dm+\frac{3d-d^2}{2}.
\]
Porównując wyraz stały z zapisem $dm+1-p_a$,
otrzymujemy $p_a=1-(3d-d^2)/2=(d-1)(d-2)/2$.
\end{proof}

\begin{uwaga}[Rodzaj topologiczny może być mniejszy]
Dla gładkiej krzywej płaskiej nad $\C$ rodzaj
arytmetyczny zgadza się z topologicznym rodzajem
powierzchni Riemanna; dowód tej głębokiej zgodności
przeprowadzimy dopiero po wprowadzeniu dywizorów
i twierdzenia Riemanna--Roch. Dla osobliwej krzywej
kubicznej $p_a=1$, chociaż normalizacje zarówno
kubiki z ostrzem, jak i z węzłem są izomorficzne
z $\PP^1$, której rodzaj topologiczny jest równy $0$.
Ta rozbieżność jest informacją o osobliwości, nie błędem
w rachunku rodzaju.
\end{uwaga}

\section{Podsumowanie: trzy poziomy badania osobliwości}

Przestrzeń styczna $T_pX$ wykrywa brak zwykłego kierunku
stycznego, lecz dla dwóch różnych osobliwości może być
jednakowo duża. Stożek styczny odróżnia pojedynczą
wielokrotną styczną od kilku różnych kierunków.
Lokalny pierścień zachowuje krotności przecięć, a jego
normalizacja może rozdzielić gałęzie przechodzące przez
ten sam punkt. Rozdmuchanie zastępuje punkt przestrzenią
kierunków i pozwala geometrycznie badać powstawanie
osobliwości oraz ich rozwiązywanie.

W kolejnym rozdziale naturalnie przejdziemy do
całościowej teorii krzywych algebraicznych:
dywizorów, funkcji wymiernych, różniczek regularnych,
twierdzenia Riemanna--Roch i rodzaju krzywej.


% Rozdział 4. Algebra przemienna (wraz z going up i płaskością), snopy i krzywe.
% Numeracja zachowuje dotychczasowy układ rozdziałów 0--3.
